\chapter{Verso l'NLP: Word2Vec}
\label{cap:word2vec}
\fonti{\rif{2026-11-23}{23/11}}{7-Word2Vec, pp.\ 1--127}

Con questo capitolo inizia la parte del corso tenuta dal Prof.~Boffa. Il filo è quello del cap.~\ref{cap:embedding}: imparare
una \textbf{rappresentazione} (un embedding) per dati categorici \textbf{senza label}, con un pretext task self-supervised. Qui la
supervisione non viene dalla ricostruzione né dalle augmentation, ma dalla \textbf{struttura relazionale} dei dati: quali oggetti
compaiono vicini. Il caso più noto è il testo (Word2Vec); il cap.~\ref{cap:darkvec} applica la stessa idea al traffico di rete
(DarkVec). Il riferimento del docente per il collegamento tra pretext task e log di rete è un suo lavoro, Boffa et al., \emph{On using
pretext tasks to learn representations from network logs} (2022): Word2Vec applicato alle sessioni bash raccolte da honeypot fa emergere
gruppi di URL, percorsi e IP che compaiono spesso insieme (\rif{2026-11-23}{23/11}). Su Word2Vec non è previsto un lab.

\begin{esame}[come chiede Word2Vec il docente (\rif{2026-11-23}{23/11})]
Nessuna formula da scrivere a memoria: può dare la formula (obiettivo dello skip-gram o probabilità con la softmax) e chiedere da dove
viene, cosa rappresenta ogni termine e qual è l'intuizione. E non confondere \textbf{non contestuale} (Word2Vec, un vettore per parola)
con \textbf{contestuale} (transformer).
\end{esame}

% =====================================================================
\section{Ripasso: cosa sappiamo già degli embedding}
\begin{definizione}[embedding --- slide 7, p.~5]
Rappresentazione numerica dei dati in uno spazio vettoriale continuo in cui oggetti semanticamente o strutturalmente simili finiscono
in punti vicini. Così dati complessi si possono confrontare, raggruppare ed elaborare con operazioni matematiche.
\end{definizione}
In breve (slide 7, p.~18):
\begin{itemize}
  \item un embedding è un vettore multidimensionale che mappa proprietà del mondo reale in uno spazio numerico; è utile soprattutto per
        le \textbf{feature categoriche};
  \item è la rete a imparare come fare questa mappatura in modo significativo;
  \item \textbf{supervised}: la rete impara gli embedding risolvendo un compito con label esterne (classificazione, regressione, ma
        anche reinforcement learning, dove la ``label'' è la ricompensa);
  \item \textbf{self-supervised}: la rete impara gli embedding risolvendo un \textbf{pretext task} ricavato dai dati stessi, senza label
        esterne (contrastive loss, ricostruzione, \ldots).
\end{itemize}
La slide riporta un post di Yann LeCun del 2022 che mette il self-supervised learning al primo posto tra le idee più importanti
dell'AI dell'ultimo decennio (seguono ResNet, gating e attenzione, memoria differenziabile e moduli permutation-equivariant come la
multi-head self-attention dei transformer).

% =====================================================================
\section{Il tema: imparare dalla struttura relazionale}
Con le immagini il pretext task è intuitivo: un encoder CNN comprime l'immagine in un embedding e un decoder CNN la ricostruisce; si
impara una \textbf{compressione con perdita} (cap.~\ref{cap:embedding}). Con il testo la domanda cambia (slide 7, pp.~29--38):
\begin{itemize}
  \item nella frase ``\en{The king and the queen entered the castle}'', come si impara un embedding per le parole \emph{king} e
        \emph{queen}? E uno per l'intera frase?
  \item e per una \textbf{sequenza di IP che fanno scansione}? Si vorrebbe un embedding in cui gli IP della botnet Mirai siano lontani da
        quelli di Stretchoid.
\end{itemize}
Le risposte sono due: l'\textbf{NLP} (Word2Vec, questo capitolo) e le sue \textbf{applicazioni alla cybersecurity} (DarkVec,
cap.~\ref{cap:darkvec}).

% =====================================================================
\section{Natural Language Processing}
\begin{definizione}[NLP --- slide 7, p.~42]
Ramo dell'intelligenza artificiale che permette ai computer di comprendere, interpretare e generare il linguaggio umano.
\end{definizione}
\begin{itemize}
  \item Le parole sono \textbf{feature categoriche nominali}: il problema è lo stesso del cap.~\ref{cap:embedding}, con un vocabolario
        enorme.
  \item Nonostante gli sforzi dei linguisti, è difficile scrivere regole sistematiche che mappino le parole in vettori: intuitivamente
        \emph{king} e \emph{queen} dovrebbero avere embedding simili, e \emph{king} e \emph{kangaroo}?
  \item Applicazioni note: sentiment analysis, categorizzazione di testi, traduzione automatica, question answering, topic modelling,
        riassunto automatico.
  \item Oggi la modalità testuale funziona così bene che si tende a \textbf{trasformare tutto in testo}: anche immagini e grafi vengono
        affrontati con modelli di linguaggio (la slide cita due articoli, uno per le immagini e uno per i grafi).
\end{itemize}

% =====================================================================
\section{Prima delle reti neurali: bag of words}
\textbf{Idea}: ogni frase è un ``sacchetto'' di parole. Per ogni parola del vocabolario si scrive quante volte compare nella frase;
le parole assenti valgono 0. La frase diventa un vettore lungo quanto il vocabolario (slide 7, pp.~49--54).
\begin{itemize}
  \item Il \textbf{vocabolario} può essere noto (es.\ un dizionario inglese) oppure \textbf{stimato} dai dati. In quel caso, su quali
        dati? Solo sul training, come per qualsiasi vocabolario (cap.~\ref{cap:embedding}).
  \item Il docente la chiama \en{frequency-based baseline}; il nome della libreria è \textbf{count vectorizer}
        (\code{sklearn.feature\_extraction}). Una variante pesa i termini per scoraggiare quelli troppo frequenti: \textbf{TF-IDF}.
  \item È la stessa rappresentazione usata come baseline nel Lab~2 sulle sequenze di API call (cap.~\ref{cap:lab2}).
\end{itemize}
\begin{attenzione}[una baseline solida ma cieca]
La bag of words funziona spesso bene, a volte in modo frustrante per chi prova modelli più complessi. Però perde del tutto
l'\textbf{ordine} delle parole e la \textbf{semantica} (cosa significano): guarda solo come è scritta la frase. Con un vocabolario
grande soffre la \textbf{maledizione della dimensionalità}.
\end{attenzione}

% =====================================================================
\section{Word2Vec}
Word2Vec (2013) è un algoritmo per estrarre embedding di parole. Si basa sull'\textbf{ipotesi distribuzionale}, un esempio di
\en{model engineering} (si progetta il modello attorno a un'ipotesi sui dati):
\begin{definizione}[ipotesi distribuzionale --- slide 7, pp.~58--62]
Il significato di una parola dipende dal \textbf{contesto} in cui compare. Non da tutto il contesto, ma soprattutto dalle parole
\textbf{più vicine}: è l'intuizione chiave di Word2Vec.
\end{definizione}
A lezione il docente la riconduce al linguista J.~R.~Firth (1957): ``\en{you shall know a word by the company it keeps}''. Più esempi
di una parola si hanno, più il suo significato è affidabile. In termini di rete: se dall'embedding di \emph{banking} si riescono a
predire le parole che lo circondano (o viceversa), l'embedding ha catturato qualcosa del significato.
Esempio della slide: ``\ldots government debt problems turning into \textbf{banking} crises as has happened in\ldots'' e ``\ldots saying
that Europe handles the \textbf{banking} crisis with a shared market\ldots'': le parole attorno a \emph{banking} ne definiscono il senso.

Due modi di trasformare l'ipotesi in un compito per una rete:
\begin{center}
\small
\begin{tabularx}{\textwidth}{l X}
\toprule
\textbf{Variante} & \textbf{Compito}\\
\midrule
CBOW (\en{Continuous Bag of Words}) & predire la parola centrale date le parole del contesto\\
\textbf{Skip-gram} & predire le parole del contesto data la parola centrale (la più usata, e quella trattata dal corso)\\
\bottomrule
\end{tabularx}
\end{center}
Sono due articoli diversi degli stessi autori; nelle librerie la scelta tra le due varianti è un parametro.

\subsection{Lo skip-gram: l'obiettivo}
Per ogni parola $w_t$ di una frase di lunghezza $T$ si massimizza la probabilità delle parole che le stanno attorno, entro un
\textbf{contesto} $c$ (iperparametro: un contesto più largo porta forse più informazione, di sicuro più calcolo). Si massimizza la
log-probabilità media (slide 7, p.~70):
\[
  \frac{1}{T}\sum_{t=1}^{T}\ \sum_{\substack{-c \le j \le c\\ j \neq 0}} \log p(w_{t+j} \mid w_t)
\]
Letta da dentro a fuori: la probabilità di una parola di contesto data quella centrale, per tutte le parole del contesto
$[-c, c]$, per tutte le parole della frase, per tutte le frasi del corpus. Di solito $c$ si sceglie piccolo, per il costo.

\subsection{Come si stima \texorpdfstring{$p(w_{t+j} \mid w_t)$}{p(w\_t+j | w\_t)}}
\textbf{Intuizione}: data la parola centrale, il modello deve capire che le parole del suo contesto sono \textbf{più probabili} di tutte
le altre del vocabolario. Formalmente Word2Vec è una rete con \textbf{due layer lineari}, cioè due tabelle di embedding (slide 7,
pp.~81--86):
\begin{itemize}
  \item una per le parole \textbf{centrali} ($e_t$);
  \item una per le parole di \textbf{contesto} ($e_{t+j}$, e $e_w$ per una parola qualsiasi del vocabolario).
\end{itemize}
Le dimensioni aiutano (\rif{2026-11-23}{23/11}): il one-hot della parola ($1 \times V$) per una matrice $V \times H$ dà un embedding
$1 \times H$, con $H \ll V$. Lo spazio più piccolo è una riduzione di dimensionalità: costringe la rete a tenere solo le dimensioni
importanti. Prima del training le matrici sono casuali e le proiezioni non hanno alcun significato.
\[
  p(w_{t+j} \mid w_t) = \frac{\exp\!\left(e_{t+j}^{\top} e_t\right)}{\sum_{w=1}^{V} \exp\!\left(e_w^{\top} e_t\right)}
\]
\begin{itemize}
  \item Al numeratore lo \textbf{score} tra embedding centrale e di contesto: è un prodotto scalare, massimo quando i due vettori sono
        paralleli, cioè portano informazione simile (come due segnali con la stessa direzione). Massimizzarlo per le parole dello stesso
        contesto le rende simili alla parola centrale.
  \item Al denominatore lo score contro \textbf{tutte} le parole del vocabolario $V$.
  \item È una \textbf{softmax}: risolvere Word2Vec equivale a risolvere un \textbf{problema di classificazione} con $V$ classi
        (``quale parola compare nel contesto?''), con le label ricavate dal testo stesso.
\end{itemize}

\subsection{Negative sampling}
In pratica il denominatore è \textbf{incalcolabile}: bisognerebbe calcolare lo score con ogni parola del vocabolario a ogni passo, e
$V$ è enorme (la slide cita 171\,476 parole nel dizionario inglese). Si approssima con il \textbf{negative sampling} (slide 7, p.~93;
$v_{w_I}$ è l'embedding della parola centrale, $v'_{w_O}$ quello della parola di contesto osservata):
\[
  \log \sigma\!\left({v'_{w_O}}^{\top} v_{w_I}\right) + \sum_{i=1}^{k} \mathbb{E}_{w_i \sim P_n(w)}
  \left[\log \sigma\!\left(-{v'_{w_i}}^{\top} v_{w_I}\right)\right]
\]
\begin{itemize}
  \item \textbf{contributo positivo}: la coppia (centrale, contesto) vera deve avere score alto;
  \item \textbf{contributi negativi}: $k$ parole estratte a caso dal vocabolario (i \en{negative}, $k \ll V$) devono avere score basso;
  \item $k$ più grande = approssimazione migliore, ma più calcolo. Nell'articolo si mostra che l'approssimazione converge alla
        distribuzione vera al crescere di $k$.
\end{itemize}
Per il corso basta l'intuizione: chi sta nello stesso contesto deve essere vicino, chi non ci sta lontano.
È di nuovo una scelta tra coppie vicine e lontane, come la contrastive loss del cap.~\ref{cap:embedding}.

\subsection{Come si scrive}
Lo scheletro in PyTorch mostrato a lezione (slide 7, p.~98) ha \textbf{due \code{nn.Embedding}} (ricorda: un embedding layer
equivale a one-hot più un layer lineare, cap.~\ref{cap:embedding}) e una loss con i due contributi; senza negativi lo score andrebbe
calcolato su tutto il vocabolario. Qui una versione ridotta:
\begin{lstlisting}
class Word2VecSkipGram(nn.Module):
    def __init__(self, vocab_size, embed_dim):
        super().__init__()
        self.center_embeddings = nn.Embedding(vocab_size, embed_dim)   # parole centrali
        self.context_embeddings = nn.Embedding(vocab_size, embed_dim)  # parole di contesto

    def forward(self, center, context, negatives):
        # center: [B], context: [B, C], negatives: [B, C, K]
        e_c = self.center_embeddings(center).unsqueeze(1)              # [B, 1, H]
        pos = (e_c * self.context_embeddings(context)).sum(-1)        # [B, C]
        loss = -torch.log(torch.sigmoid(pos) + 1e-10).mean()          # contributo positivo
        neg = (e_c.unsqueeze(2) * self.context_embeddings(negatives)).sum(-1)  # [B, C, K]
        loss += -torch.log(torch.sigmoid(-neg) + 1e-10).mean()        # contributi negativi
        return loss
\end{lstlisting}
Nella pratica si usa una libreria: con \textbf{Gensim} basta una riga, \code{Word2Vec(sentences=sentences)} dal modulo
\code{gensim.models}, con parametri per la finestra di contesto, la variante e così via.

% =====================================================================
\section{Modelli pre-addestrati}
Non serve addestrare un modello nuovo ogni volta (slide 7, p.~103): il training self-supervised impara embedding \textbf{generici};
il pretext task si fa \textbf{una volta}, poi si fa fine-tuning. Esempio: un servizio di completamento automatico in inglese.
\begin{enumerate}
  \item si carica un Word2Vec pre-addestrato (es.\ su Wikipedia);
  \item lo si usa così com'è, come dizionario parola $\to$ embedding, oppure lo si \textbf{rifinisce} (fine-tuning) per esempio per
        preferire certi suggerimenti.
\end{enumerate}
Il pre-training non richiede label, solo frasi: ogni frase di Wikipedia è un esempio. Sopra gli embedding si può anche costruire un
classificatore: in DarkVec, per esempio, decidere se l'embedding di un IP appartiene a uno scanner benigno o alla botnet Mirai. La
stessa logica (pre-training self-supervised + fine-tuning) tornerà con i transformer.

% =====================================================================
\section{I limiti di Word2Vec}
\begin{enumerate}
  \item \textbf{Embedding non contestuali.} Ogni parola ha un solo vettore. Se si addestra su tre corpora con lo stesso numero di frasi
        (biologia e zoologia, favole per bambini, libri di informatica), dove finisce \emph{mouse}? Due corpora su tre parlano
        dell'animale: l'embedding si sposta verso l'animale, ma non rappresenta bene \textbf{nessuno} dei due significati. In realtà è
        il contesto della frase a definire la parola (``il topolino della storia ha battuto il gatto'' contro ``Engelbart ha inventato
        il mouse nel 1968''). Per adattarlo servirebbe riaddestrare su nuovi corpora.
  \item \textbf{Parole nuove.} Tutto parte da un one-hot su un vocabolario fisso: una parola nuova è \textbf{fuori vocabolario} (OOV).
        Esempio: ``petaloso'', nata nel 2016, non c'è nel Word2Vec italiano del 2015. Aggiungere una parola vuol dire allargare la
        matrice $V \times H$ di una riga e riaddestrare. Lo stesso vale per DarkVec: che fare con un IP mai visto
        (cap.~\ref{cap:darkvec})?
  \item \textbf{Parole rare.} In tedesco si compongono parole per esprimere concetti nuovi (\emph{Sitzpinkler} da \emph{sitzen} e
        \emph{pinkeln}). Word2Vec vede la parola composta e le sue parti come tre voci separate del dizionario: la composta è rara, quindi
        è difficile imparare un buon embedding. Per una parola rara di cybersecurity, o un IP raro, sarebbe invece importante. I
        modelli successivi lavorano a \textbf{token}, pezzi di parola (i limiti dei chatbot si contano in token): la parola composta
        si spezza in parti già note.
  \item \textbf{Dalle parole alle frasi.} I compiti interessanti sono quasi sempre su frasi (sentiment analysis di una recensione,
        traduzione). Con Word2Vec non è banale passare da molti oggetti (le parole) a una sola predizione (recensione positiva): serve
        un'aggregazione, e Word2Vec non dice come impararla. In sicurezza interessa per esempio decidere se un intero frammento di codice
        è malware o goodware.
\end{enumerate}
Questi limiti preparano i modelli successivi (transformer), trattati nelle lezioni seguenti.

\begin{sintesi}
\begin{itemize}
  \item Le parole sono feature categoriche nominali con vocabolari enormi: servono embedding, e non ci sono label.
  \item \textbf{Bag of words} (count vectorizer, TF-IDF): baseline solida, ma perde ordine e semantica e cresce con il vocabolario.
  \item \textbf{Word2Vec}: ipotesi distribuzionale (il significato viene dalle parole vicine); \textbf{CBOW} predice la parola
        centrale, \textbf{skip-gram} predice il contesto. Due tabelle di embedding; score = prodotto scalare; softmax su $V$ classi
        $\Rightarrow$ è una classificazione con label ricavate dal testo (self-supervised).
  \item \textbf{Negative sampling}: un positivo e $k$ negativi casuali al posto della softmax completa.
  \item Pre-addestrato una volta, poi riusato o rifinito.
  \item Limiti: embedding non contestuali, OOV, parole rare, aggregazione da parole a frasi.
\end{itemize}
\end{sintesi}
